% ============================================================
% 04-methodology.tex
% ============================================================
\section{Methodology}
\label{sec:method}

This section describes the core methodologies implemented in CTU-Chat,
including hybrid retrieval, cross-encoder evidence reranking, governed
retrieval lanes, graph-based candidate expansion, and domain-specific
deterministic calculation tools.

%------------------------------------------------
\subsection{Hybrid Retrieval with Reciprocal Rank Fusion}

To improve retrieval robustness, CTU-Chat combines BM25 lexical retrieval and
dense semantic retrieval. Since these retrievers produce heterogeneous score
distributions, their scores are not directly summed. Instead, retrieved
candidates are merged, deduplicated by document identifier, and ranked using
Reciprocal Rank Fusion (RRF). Graph-derived candidates are added later by the
orchestration layer and are reranked with the fused text candidates.

For the retrieval source set
$\mathcal{S}=\{\text{BM25},\text{Dense}\}$,
the RRF score of a candidate document $d$ is computed as

\begin{equation}
  \RRF(d)
  =\sum_{s\in\mathcal{S}_d}\frac{1}{k+r_s(d)},
  \qquad k=60,
\label{eq:rrf}
\end{equation}

where $r_s(d)$ denotes the ranking position of document $d$ returned by retrieval source $s$, and $\mathcal{S}_d$ contains only the retrievers that successfully retrieve $d$. The constant $k=60$ is fixed across all experiments to ensure reproducibility.

The query-intent classifier selects metadata-constrained retrieval lanes for
tuition, scholarships, loans, social support, academic programs, and academic
regulations. For supported academic-program queries, Neo4j entity lookup adds
program metadata to the candidate set. The graph is treated as an expansion
mechanism rather than an authority source.

%------------------------------------------------
\subsection{Cross-Encoder Evidence Reranking}

Following RRF fusion and governed-lane expansion, CTU-Chat uses
\nolinkurl{BAAI/bge-reranker-v2-m3} to score each query--document pair. The
implementation truncates the text used for scoring to 1,000 characters while
retaining the full document for answer generation. Documents are ranked by
cross-encoder relevance; when scores differ by at most 0.05, newer timestamped
documents are preferred. This neural reranker is not the hand-weighted
deterministic formula described in earlier drafts.

%------------------------------------------------
\subsection{Deterministic Academic and Financial Tools}

Certain university services require exact numerical computation rather than language-model reasoning. CTU-Chat therefore integrates typed deterministic tools for course grading, GPA calculation, and tuition fee estimation, ensuring reproducible results that are directly traceable to institutional policies.

\subsubsection{Course Grade Calculation}

Given assessment scores $x_i$ and weights $w_i$ satisfying $\sum_i w_i=100$, the final 10-point score is calculated as

\begin{equation}
  S_{10}
  =\mathrm{Round}_{0.1}\!\left(
     \sum_i
     \mathrm{Round}_{0.1}(x_i)
     \frac{w_i}{100}
   \right).
\label{eq:grade}
\end{equation}

The resulting score is converted into the official CTU letter-grade scale (A, B+, B, C+, C, D+, D, and F).

\subsubsection{GPA Calculation}

For eligible courses $\mathcal{I}$ with credits $c_i$ and grade points $g_i$, GPA is computed as

\begin{equation}
  \mathrm{GPA}
  =\mathrm{Round}_{0.01}\!\left(
     \frac{\sum_{i\in\mathcal{I}}c_i g_i}
          {\sum_{i\in\mathcal{I}}c_i}
   \right).
\label{eq:gpa}
\end{equation}

Courses with statuses Missing (M), Incomplete (I), and Withdrawn (W) are excluded according to the university grading policy.

\subsubsection{Tuition Fee Calculation}

For each registered course $i$, the payable tuition fee is determined from the credit load, tuition rate, exemption basis, and reduction percentage:

\begin{align}
  \mathrm{Gross}_i &= c_i r_i m_i,\\
  \mathrm{Reduction}_i &= c_i b_i\frac{p_i}{100},\\
  \mathrm{Payable}_i &=
  \mathrm{Round}_{1\,\mathrm{VND}}
  \left(\max(0,\mathrm{Gross}_i-\mathrm{Reduction}_i)\right).
\label{eq:tuition-payable}
\end{align}

All computations use decimal arithmetic, and every tool output preserves the corresponding evidence identifier and policy reference to guarantee auditability.
